\documentclass[11pt]{article}
    \usepackage[utf8]{inputenc}
    \usepackage[T1]{fontenc}
    \usepackage[margin=1in]{geometry}
    \usepackage{hyperref}
    \usepackage{enumitem}
    \title{Lego Spike: Swamp Boat}
    \author{Piter Garcia}
    \date{March 20, 2026}
    \begin{document}
    \maketitle
    \section*{Quick Scan}
    \begin{itemize}[leftmargin=*]
    \item Grade: Grade 5
    \item Workbook sessions: Session 17
    \item Class blocks: Day 2 10:45-11:35 Pickett, Day 3 10:45-11:35 McGlashon, Day 5 10:45-11:35 Pecora
    \item Reference file: lego-spike-swamp-boat.bib
    \end{itemize}
    \section*{School Planning Goals and Inclusion}
    \begin{itemize}[leftmargin=*]
    \item What is expected: I can explain how an object's speed is related to its energy, using evidence. | I can create explain how to protect the environment | I can iterate and test my models to improve them.
    \item What we achieved: No direct transcript match is attached yet. Treat this lesson as workbook-backed and enrich it when a matching classroom transcript is available.
    \item What needs improvement: stronger proof, cleaner assessment notes, and tighter lesson artifacts.
    \item How to improve: add transcript proof, one saved artifact, and clearer success criteria.
    \item Inclusion focus: use visual modeling, chunked directions, predictable routines, and flexible participation roles.
    \end{itemize}
    \section*{Official References}
    \begin{itemize}[leftmargin=*]
    \item \url{https://csteachers.org/k12standards/interactive/}
\item \url{https://assets.education.lego.com/v3/assets/blt293eea581807678a/blt593e0d51d19eceb5/62f4d13ce894104b27259cb8/Coding_express_teacher_guide_110822.pdf?locale=ja-jp}

    \end{itemize}
    \end{document}
